\documentclass[11pt,a4paper]{article}
\usepackage[utf8]{inputenc}
\usepackage[T1]{fontenc}
\usepackage[spanish,es-minimal,es-noshorthands]{babel}
\AtBeginDocument{\renewcommand{\sin}{\operatorname{sen}}\renewcommand{\arcsin}{\operatorname{arcsen}}}
\usepackage{lmodern}
\usepackage[sfdefault,lining]{FiraSans}
\usepackage[a4paper,top=2.3cm,bottom=2.2cm,left=1.8cm,right=1.8cm,headheight=1.3cm,headsep=0.4cm,footskip=1cm]{geometry}
\usepackage{amsmath,amssymb}
\usepackage{icomma}
\usepackage{graphicx}
\usepackage[table]{xcolor}
\usepackage{tikz}
\usepackage{fontawesome5}
\usepackage{ccicons}
\usepackage{fancyhdr}
\usepackage{lastpage}
\usepackage{tasks}
\usepackage{needspace}
\usepackage[most]{tcolorbox}
\usepackage{array}
\usepackage[hidelinks]{hyperref}
\AtBeginDocument{\hypersetup{pdftitle={Ficha de ejercicios de la unidad 10: Distribuciones de probabilidad. 1º Bach.},pdfsubject={Matemáticas CCSS 1º Bach.},pdfkeywords={pdfkeywords={variable aleatoria, función de probabilidad, función de distribución, esperanza, distribución binomial, función de densidad, distribución uniforme, distribución normal, normal tipificada, tipificación, aproximación de la binomial a la normal}}}}

\definecolor{navy}{HTML}{1B2748}
\definecolor{accent}{HTML}{E8603A}
\definecolor{teal}{HTML}{1F9E90}
\definecolor{soft}{HTML}{EEF1F8}
\definecolor{muted}{HTML}{5D6685}

\setlength{\parindent}{0pt}
\setlength{\parskip}{3pt}
\settasks{label=\alph*), label-format=\bfseries\color{navy}, label-width=1.3em, item-indent=1.8em, column-sep=0.6em, after-item-skip=5pt, before-skip=3pt, after-skip=4pt}

% ---------- capçalera i peu ----------
\newcommand{\llicencia}{\href{https://creativecommons.org/licenses/by-nc-sa/4.0/deed.es}{\ccbyncsa}\ \ Licencia Creative Commons Atribución-NoComercial-CompartirIgual 4.0 Internacional (CC BY-NC-SA 4.0)}
\pagestyle{fancy}
\fancyhf{}
\fancyhead[L]{\raisebox{-0.35\height}{\includegraphics[height=0.95cm]{logo-gva.png}}}
\fancyhead[C]{\small\color{muted}Ficha de ejercicios · Unidad 10: Distribuciones de probabilidad · 1\textsuperscript{o} Bach.}
\fancyhead[R]{\raisebox{-0.35\height}{\includegraphics[height=0.95cm]{logo-vermellar.png}}}
\fancyfoot[L]{\scriptsize\color{muted}\llicencia}
\fancyfoot[R]{\small\color{muted}\thepage}
\renewcommand{\headrulewidth}{0.4pt}
\renewcommand{\headrule}{\color{muted!40}\hrule height \headrulewidth}
\fancypagestyle{primera}{\fancyhf{}\renewcommand{\headrulewidth}{0pt}
  \fancyfoot[L]{\scriptsize\color{muted}\llicencia}
  \fancyfoot[R]{\small\color{muted}\thepage}}

% ---------- nivells i exercicis ----------
\newcommand{\estrelles}[1]{{\color{accent}\ifcase#1\or\faStar\faStar[regular]\faStar[regular]\or\faStar\faStar\faStar[regular]\or\faStar\faStar\faStar\fi}}
\newcounter{exer}
\newcommand{\ex}[1]{\par\needspace{4\baselineskip}\medskip\refstepcounter{exer}\noindent{\color{navy}\bfseries\theexer.}\ \ }

\newcommand{\seccio}[2]{\par\needspace{16\baselineskip}\bigskip
  \begin{tcolorbox}[colback=navy,colframe=navy,coltext=white,arc=2mm,boxrule=0pt,left=3mm,right=3mm,top=1.2mm,bottom=1.2mm]
  {\large\bfseries #1\quad #2}
  \end{tcolorbox}\vspace{-2pt}}

\usepackage{enumitem}
\newcommand{\exf}[1]{\par\medskip\refstepcounter{exer}\noindent{\color{navy}\bfseries\theexer.}\ \ }
\providecommand{\gt}{>}\providecommand{\lt}{<}
% quadre de recordatori (només a la fitxa d'exercicis)
\newcommand{\recorda}[1]{\par\smallskip\begin{tcolorbox}[colback=accent!7,colframe=accent!55,arc=2mm,boxrule=0.5pt,left=2.5mm,right=2.5mm,top=1mm,bottom=1mm]\small{\color{accent}\faLightbulb[regular]}\ \textbf{Recuerda:}\ #1\end{tcolorbox}}
\newcommand{\pau}[1]{{\color{teal}\small\bfseries\faUniversity\ PAU · #1.}\ }
\newcommand{\res}[1]{#1}\newcommand{\rest}[1]{#1}

\def\fmneed{11}% que el títol d'una secció no quede sol al final del full
% =====================================================================
% Capçalera FentMates de les fitxes: un color per a cada curs.
% Els colors són els de la pàgina inicial (LEVELS a assets/js/materials.js);
% el logo ve de _eines/logo/pdf/ (python3 _eines/logo/genera.py).
%
% Ús, al final del preàmbul de la unitat:
%   % =====================================================================
% Capçalera FentMates de les fitxes: un color per a cada curs.
% Els colors són els de la pàgina inicial (LEVELS a assets/js/materials.js);
% el logo ve de _eines/logo/pdf/ (python3 _eines/logo/genera.py).
%
% Ús, al final del preàmbul de la unitat:
%   % =====================================================================
% Capçalera FentMates de les fitxes: un color per a cada curs.
% Els colors són els de la pàgina inicial (LEVELS a assets/js/materials.js);
% el logo ve de _eines/logo/pdf/ (python3 _eines/logo/genera.py).
%
% Ús, al final del preàmbul de la unitat:
%   \input{../comu/fentmates.tex}
%   \fmcurs{1eso}
% i, al començament del document:
%   \capcalera{Matemàtiques}{1r ESO}{Fitxa d'exercicis · Unitat 2}{Divisibilitat}
% Opcional abans de \input: \def\fmneed{14} (espai mínim abans de cada secció).
% =====================================================================
\makeatletter
\def\fm@logodir{../../logo/pdf}
\providecommand{\fmneed}{10}
\def\fm@cursnom{}

% clau del curs, color i color del text damunt del color
\newcommand{\fm@curs}[3]{\@namedef{fm@color@#1}{#2}\@namedef{fm@text@#1}{#3}}
\fm@curs{1eso}{E6B84A}{navy}
\fm@curs{2eso}{3B82F6}{white}
\fm@curs{3eso}{FF7A50}{white}
\fm@curs{4eso-a}{A78BFA}{white}
\fm@curs{4eso-b}{6D28D9}{white}
\fm@curs{1batx-mi}{9B1C1C}{white}
\fm@curs{1batx-ccss}{EF6B6B}{white}
\fm@curs{2batx-mii}{1E6B47}{white}
\fm@curs{2batx-ccss}{2BB3A3}{white}

\newcommand{\fmcurs}[1]{%
  \@ifundefined{fm@color@#1}{\PackageError{fentmates}{Curs desconegut: #1}{}}{}%
  \edef\fm@tmp{\noexpand\definecolor{curs}{HTML}{\@nameuse{fm@color@#1}}%
               \noexpand\colorlet{curstext}{\@nameuse{fm@text@#1}}}\fm@tmp
  \colorlet{cursfosc}{curs!72!black}%
  \def\fm@logo{\fm@logodir/vertical-#1.pdf}%
  \def\fm@marca{\fm@logodir/marca-#1.pdf}}

% ---------- capçalera de la primera pàgina ----------
% #1 matèria (petit), #2 curs, #3 tipus de document i unitat, #4 títol
\newcommand{\capcalera}[4]{%
  \gdef\fm@cursnom{#2}%
  \par\noindent
  \begin{tcolorbox}[colback=curs,colframe=curs,boxrule=0pt,arc=2.2mm,width=3.5cm,height=2.15cm,
      valign=center,left=3mm,right=2mm,top=1mm,bottom=1mm,nobeforeafter,box align=center]
    \color{curstext}{\footnotesize #1}\par\vspace{1pt}{\fontsize{19}{21}\selectfont\bfseries #2\par}%
  \end{tcolorbox}\hspace{5mm}%
  \begin{minipage}[c]{\dimexpr\linewidth-3.5cm-5mm-2.9cm\relax}
    {\color{cursfosc}\small\bfseries\MakeUppercase{#3}}\par\vspace{2pt}
    {\color{navy}\fontsize{20}{23}\selectfont\bfseries #4\par}
  \end{minipage}\hfill
  \begin{minipage}[c]{2.6cm}\raggedleft\includegraphics[width=2.6cm]{\fm@logo}\end{minipage}\par
  \vspace{2pt}}

% ---------- barres de secció del color del curs ----------
\renewcommand{\seccio}[2]{\par\needspace{\fmneed\baselineskip}\bigskip
  \begin{tcolorbox}[colback=curs,colframe=curs,coltext=curstext,arc=2mm,boxrule=0pt,left=3mm,right=3mm,top=1.2mm,bottom=1.2mm]
  {\large\bfseries #1\quad #2}
  \end{tcolorbox}\vspace{-2pt}}

% ---------- estrelles de dificultat del color del curs (a Batxillerat no se'n posen) ----------
\@ifundefined{estrelles}{}{%
  \renewcommand{\estrelles}[1]{{\color{curs}\ifcase#1\or\faStar\faStar[regular]\faStar[regular]\or\faStar\faStar\faStar[regular]\or\faStar\faStar\faStar\fi}}}

% ---------- capçalera i peu de les pàgines ----------
% primera pàgina: logos institucionals al peu
\fancypagestyle{primera}{\fancyhf{}\renewcommand{\headrulewidth}{0pt}
  \fancyfoot[L]{\raisebox{-0.32\height}{\includegraphics[height=0.78cm]{logo-gva.png}}\hspace{5mm}%
    \raisebox{-0.32\height}{\includegraphics[height=0.78cm]{logo-vermellar.png}}}
  \fancyfoot[R]{\scriptsize\color{muted}\href{https://creativecommons.org/licenses/by-nc-sa/4.0/deed.ca}{\ccbyncsa}\ \ CC BY-NC-SA 4.0 · fentmates.org\qquad\small\thepage}}
% la resta: el símbol a l'esquerra i el curs a la dreta, amb el títol al centre com abans
\fancyhead[L]{\raisebox{-0.3\height}{\includegraphics[height=0.62cm]{\fm@marca}}}
\fancyhead[R]{\small\bfseries\color{cursfosc}\fm@cursnom}
\renewcommand{\headrule}{\color{curs!45}\hrule height \headrulewidth}
\makeatother

%   \fmcurs{1eso}
% i, al començament del document:
%   \capcalera{Matemàtiques}{1r ESO}{Fitxa d'exercicis · Unitat 2}{Divisibilitat}
% Opcional abans de \input: \def\fmneed{14} (espai mínim abans de cada secció).
% =====================================================================
\makeatletter
\def\fm@logodir{../../logo/pdf}
\providecommand{\fmneed}{10}
\def\fm@cursnom{}

% clau del curs, color i color del text damunt del color
\newcommand{\fm@curs}[3]{\@namedef{fm@color@#1}{#2}\@namedef{fm@text@#1}{#3}}
\fm@curs{1eso}{E6B84A}{navy}
\fm@curs{2eso}{3B82F6}{white}
\fm@curs{3eso}{FF7A50}{white}
\fm@curs{4eso-a}{A78BFA}{white}
\fm@curs{4eso-b}{6D28D9}{white}
\fm@curs{1batx-mi}{9B1C1C}{white}
\fm@curs{1batx-ccss}{EF6B6B}{white}
\fm@curs{2batx-mii}{1E6B47}{white}
\fm@curs{2batx-ccss}{2BB3A3}{white}

\newcommand{\fmcurs}[1]{%
  \@ifundefined{fm@color@#1}{\PackageError{fentmates}{Curs desconegut: #1}{}}{}%
  \edef\fm@tmp{\noexpand\definecolor{curs}{HTML}{\@nameuse{fm@color@#1}}%
               \noexpand\colorlet{curstext}{\@nameuse{fm@text@#1}}}\fm@tmp
  \colorlet{cursfosc}{curs!72!black}%
  \def\fm@logo{\fm@logodir/vertical-#1.pdf}%
  \def\fm@marca{\fm@logodir/marca-#1.pdf}}

% ---------- capçalera de la primera pàgina ----------
% #1 matèria (petit), #2 curs, #3 tipus de document i unitat, #4 títol
\newcommand{\capcalera}[4]{%
  \gdef\fm@cursnom{#2}%
  \par\noindent
  \begin{tcolorbox}[colback=curs,colframe=curs,boxrule=0pt,arc=2.2mm,width=3.5cm,height=2.15cm,
      valign=center,left=3mm,right=2mm,top=1mm,bottom=1mm,nobeforeafter,box align=center]
    \color{curstext}{\footnotesize #1}\par\vspace{1pt}{\fontsize{19}{21}\selectfont\bfseries #2\par}%
  \end{tcolorbox}\hspace{5mm}%
  \begin{minipage}[c]{\dimexpr\linewidth-3.5cm-5mm-2.9cm\relax}
    {\color{cursfosc}\small\bfseries\MakeUppercase{#3}}\par\vspace{2pt}
    {\color{navy}\fontsize{20}{23}\selectfont\bfseries #4\par}
  \end{minipage}\hfill
  \begin{minipage}[c]{2.6cm}\raggedleft\includegraphics[width=2.6cm]{\fm@logo}\end{minipage}\par
  \vspace{2pt}}

% ---------- barres de secció del color del curs ----------
\renewcommand{\seccio}[2]{\par\needspace{\fmneed\baselineskip}\bigskip
  \begin{tcolorbox}[colback=curs,colframe=curs,coltext=curstext,arc=2mm,boxrule=0pt,left=3mm,right=3mm,top=1.2mm,bottom=1.2mm]
  {\large\bfseries #1\quad #2}
  \end{tcolorbox}\vspace{-2pt}}

% ---------- estrelles de dificultat del color del curs (a Batxillerat no se'n posen) ----------
\@ifundefined{estrelles}{}{%
  \renewcommand{\estrelles}[1]{{\color{curs}\ifcase#1\or\faStar\faStar[regular]\faStar[regular]\or\faStar\faStar\faStar[regular]\or\faStar\faStar\faStar\fi}}}

% ---------- capçalera i peu de les pàgines ----------
% primera pàgina: logos institucionals al peu
\fancypagestyle{primera}{\fancyhf{}\renewcommand{\headrulewidth}{0pt}
  \fancyfoot[L]{\raisebox{-0.32\height}{\includegraphics[height=0.78cm]{logo-gva.png}}\hspace{5mm}%
    \raisebox{-0.32\height}{\includegraphics[height=0.78cm]{logo-vermellar.png}}}
  \fancyfoot[R]{\scriptsize\color{muted}\href{https://creativecommons.org/licenses/by-nc-sa/4.0/deed.ca}{\ccbyncsa}\ \ CC BY-NC-SA 4.0 · fentmates.org\qquad\small\thepage}}
% la resta: el símbol a l'esquerra i el curs a la dreta, amb el títol al centre com abans
\fancyhead[L]{\raisebox{-0.3\height}{\includegraphics[height=0.62cm]{\fm@marca}}}
\fancyhead[R]{\small\bfseries\color{cursfosc}\fm@cursnom}
\renewcommand{\headrule}{\color{curs!45}\hrule height \headrulewidth}
\makeatother

%   \fmcurs{1eso}
% i, al començament del document:
%   \capcalera{Matemàtiques}{1r ESO}{Fitxa d'exercicis · Unitat 2}{Divisibilitat}
% Opcional abans de \input: \def\fmneed{14} (espai mínim abans de cada secció).
% =====================================================================
\makeatletter
\def\fm@logodir{../../logo/pdf}
\providecommand{\fmneed}{10}
\def\fm@cursnom{}

% clau del curs, color i color del text damunt del color
\newcommand{\fm@curs}[3]{\@namedef{fm@color@#1}{#2}\@namedef{fm@text@#1}{#3}}
\fm@curs{1eso}{E6B84A}{navy}
\fm@curs{2eso}{3B82F6}{white}
\fm@curs{3eso}{FF7A50}{white}
\fm@curs{4eso-a}{A78BFA}{white}
\fm@curs{4eso-b}{6D28D9}{white}
\fm@curs{1batx-mi}{9B1C1C}{white}
\fm@curs{1batx-ccss}{EF6B6B}{white}
\fm@curs{2batx-mii}{1E6B47}{white}
\fm@curs{2batx-ccss}{2BB3A3}{white}

\newcommand{\fmcurs}[1]{%
  \@ifundefined{fm@color@#1}{\PackageError{fentmates}{Curs desconegut: #1}{}}{}%
  \edef\fm@tmp{\noexpand\definecolor{curs}{HTML}{\@nameuse{fm@color@#1}}%
               \noexpand\colorlet{curstext}{\@nameuse{fm@text@#1}}}\fm@tmp
  \colorlet{cursfosc}{curs!72!black}%
  \def\fm@logo{\fm@logodir/vertical-#1.pdf}%
  \def\fm@marca{\fm@logodir/marca-#1.pdf}}

% ---------- capçalera de la primera pàgina ----------
% #1 matèria (petit), #2 curs, #3 tipus de document i unitat, #4 títol
\newcommand{\capcalera}[4]{%
  \gdef\fm@cursnom{#2}%
  \par\noindent
  \begin{tcolorbox}[colback=curs,colframe=curs,boxrule=0pt,arc=2.2mm,width=3.5cm,height=2.15cm,
      valign=center,left=3mm,right=2mm,top=1mm,bottom=1mm,nobeforeafter,box align=center]
    \color{curstext}{\footnotesize #1}\par\vspace{1pt}{\fontsize{19}{21}\selectfont\bfseries #2\par}%
  \end{tcolorbox}\hspace{5mm}%
  \begin{minipage}[c]{\dimexpr\linewidth-3.5cm-5mm-2.9cm\relax}
    {\color{cursfosc}\small\bfseries\MakeUppercase{#3}}\par\vspace{2pt}
    {\color{navy}\fontsize{20}{23}\selectfont\bfseries #4\par}
  \end{minipage}\hfill
  \begin{minipage}[c]{2.6cm}\raggedleft\includegraphics[width=2.6cm]{\fm@logo}\end{minipage}\par
  \vspace{2pt}}

% ---------- barres de secció del color del curs ----------
\renewcommand{\seccio}[2]{\par\needspace{\fmneed\baselineskip}\bigskip
  \begin{tcolorbox}[colback=curs,colframe=curs,coltext=curstext,arc=2mm,boxrule=0pt,left=3mm,right=3mm,top=1.2mm,bottom=1.2mm]
  {\large\bfseries #1\quad #2}
  \end{tcolorbox}\vspace{-2pt}}

% ---------- estrelles de dificultat del color del curs (a Batxillerat no se'n posen) ----------
\@ifundefined{estrelles}{}{%
  \renewcommand{\estrelles}[1]{{\color{curs}\ifcase#1\or\faStar\faStar[regular]\faStar[regular]\or\faStar\faStar\faStar[regular]\or\faStar\faStar\faStar\fi}}}

% ---------- capçalera i peu de les pàgines ----------
% primera pàgina: logos institucionals al peu
\fancypagestyle{primera}{\fancyhf{}\renewcommand{\headrulewidth}{0pt}
  \fancyfoot[L]{\raisebox{-0.32\height}{\includegraphics[height=0.78cm]{logo-gva.png}}\hspace{5mm}%
    \raisebox{-0.32\height}{\includegraphics[height=0.78cm]{logo-vermellar.png}}}
  \fancyfoot[R]{\scriptsize\color{muted}\href{https://creativecommons.org/licenses/by-nc-sa/4.0/deed.ca}{\ccbyncsa}\ \ CC BY-NC-SA 4.0 · fentmates.org\qquad\small\thepage}}
% la resta: el símbol a l'esquerra i el curs a la dreta, amb el títol al centre com abans
\fancyhead[L]{\raisebox{-0.3\height}{\includegraphics[height=0.62cm]{\fm@marca}}}
\fancyhead[R]{\small\bfseries\color{cursfosc}\fm@cursnom}
\renewcommand{\headrule}{\color{curs!45}\hrule height \headrulewidth}
\makeatother

\fmcurs{1batx-ccss}

\begin{document}
\thispagestyle{primera}

\capcalera{Matemáticas CCSS}{1\textsuperscript{o} Bach.}{Ficha de ejercicios · Unidad 10}{\raggedright Distribuciones de probabilidad}
\vspace{10pt}
Nombre y apellidos: \rule{6.9cm}{0.4pt}\hfill Grupo: \rule{1.6cm}{0.4pt}\hfill Fecha: \rule{2.2cm}{0.4pt}
\vspace{8pt}

\begin{tcolorbox}[colback=soft,colframe=soft,arc=2mm,left=3mm,right=3mm,top=1.5mm,bottom=1.5mm]
\small Resuelve los ejercicios en el cuaderno mostrando todos los pasos: identifica la variable aleatoria y su distribución, haz un dibujo del área que buscas y redondea $z$ a dos decimales. Usa las tablas de la binomial y de la normal $N(0,1)$ que tienes en los apuntes interactivos de la unidad (son las de la PAU valenciana). Los 12 ejercicios marcados con \textcolor{teal}{\faUniversity\ PAU} son de las pruebas de acceso a la universidad.
\end{tcolorbox}

% =====================================================
\seccio{1}{Variables aleatorias discretas}

\exf{1} Di si las variables aleatorias siguientes son discretas o continuas:\begin{tasks}(2)
  \task El número de hermanos de un alumno.
  \task El tiempo que tarda un autobús en llegar.
  \task El peso de un recién nacido.
  \task El número de clientes que entran en una tienda en una hora.
  \task La temperatura máxima de un día.
  \task El número de caras en 10 lanzamientos de una moneda.
\end{tasks}

\ex{1} Lanzamos dos monedas y $X$ es el número de caras obtenidas.\begin{tasks}(1)
  \task Haz la tabla de la función de probabilidad y represéntala con un diagrama de barras.
  \task Halla la función de distribución $F(x)=P(X\leq x)$ para cada valor.
  \task Calcula la media y la desviación típica.
\end{tasks}\recorda{$\mu=\sum x_i\,p_i$; \ $\sigma^2=\sum x_i^2\,p_i-\mu^2$; \ $\sigma=\sqrt{\sigma^2}$.}

\ex{1} Una variable aleatoria $X$ toma los valores 1, 2, 3 y 4 con probabilidades $k$, $2k$, $3k$ y $4k$.\begin{tasks}(3)
  \task Calcula $k$.
  \task Calcula $\mu$ y $\sigma$.
  \task Calcula $P(X\geq3)$ y $P(X\lt3)$.
\end{tasks}

\ex{2} Lanzamos dos dados y $X$ es la suma de los puntos. Haz la tabla de la función de probabilidad y calcula $P(X\geq10)$, la media y la desviación típica.

\ex{2} Una urna contiene 4 bolas blancas y 2 negras. Sacamos dos bolas sin reemplazamiento y $X$ es el número de bolas negras. Halla la función de probabilidad de $X$, la media y la desviación típica.\recorda{haz un diagrama de árbol: la segunda extracción depende de la primera.}

\ex{2} En un juego se paga 1 euro por lanzar un dado. Si sale un 6 se cobran 4 euros; si sale un 5, 2 euros; y en caso contrario no se cobra nada. Sea $X$ la ganancia neta del jugador. Calcula la ganancia esperada y la desviación típica. ¿Es un juego justo?\recorda{la ganancia neta es lo que se cobra menos lo que se paga. Un juego es justo si la ganancia esperada es 0.}

\ex{3} Una variable aleatoria toma los valores 0, 1 y 2 con probabilidades $a$, $0{,}5$ y $b$. Si la media es $1{,}2$, calcula $a$, $b$ y la varianza.

\ex{2} La función de distribución de una variable discreta $X$ es $F(0)=0{,}1$, $F(1)=0{,}4$, $F(2)=0{,}8$ y $F(3)=1$ (y $X$ solo toma estos valores). Halla la función de probabilidad, $P(1\leq X\leq2)$ y la media.


% =====================================================
\seccio{2}{La distribución binomial}

\exf{1} Di si las variables siguientes siguen una distribución binomial y, si es así, indica sus parámetros:\begin{tasks}(1)
  \task Número de seises en 12 lanzamientos de un dado.
  \task Número de bolas rojas al sacar 3 bolas, sin reemplazamiento, de una urna con 4 rojas y 6 blancas.
  \task Número de piezas defectuosas en una muestra de 50, si el 2\,\% de la producción es defectuosa.
  \task Número de lanzamientos de una moneda hasta que sale la primera cara.
\end{tasks}\recorda{una binomial $B(n,p)$ cuenta los éxitos en $n$ pruebas independientes con la misma probabilidad de éxito $p$.}

\ex{1} Se lanza un dado 5 veces. Calcula la probabilidad de obtener exactamente dos cincos y la de obtener al menos un cinco.\recorda{$P(X=k)=\binom nk\,p^k\,q^{\,n-k}$, \ $q=1-p$.}

\ex{2} Un test tiene 8 preguntas con 3 opciones cada una, de las cuales solo una es correcta. Un alumno contesta todas las preguntas al azar.\begin{tasks}(1)
  \task ¿Qué distribución sigue el número de aciertos? ¿Cuántos aciertos se esperan? ¿Cuál es la desviación típica?
  \task Calcula la probabilidad de acertar exactamente 4 y la de acertar al menos una.
\end{tasks}

\ex{2} \pau{Castilla y León, PAU ordinaria 2026}El 90\,\% de los jóvenes de entre 25 y 34 años compra por Internet. Una tienda hace una encuesta a 10 clientes jóvenes.\begin{tasks}(1)
  \task Calcula la probabilidad de que exactamente 2 de los 10 clientes no hagan compras en línea.
  \task Calcula la probabilidad de que menos de 2 clientes no hagan compras en línea.
\end{tasks}

\ex{1} Sea $X\sim B(5;\ 0{,}25)$. Calcula con la tabla de la binomial:\begin{tasks}(4)
  \task $P(X=2)$
  \task $P(X\leq1)$
  \task $P(X\geq4)$
  \task $P(X>0)$
\end{tasks}\recorda{la tabla da $P(X=x)$ para $n\leq9$. Para las acumuladas, suma casillas de la misma columna o haz el contrario.}

\ex{2} El 80\,\% de los clientes de una tienda está satisfecho con la atención recibida. Se eligen 7 clientes al azar. Calcula, con la tabla, la probabilidad de que:\begin{tasks}(1)
  \task estén satisfechos exactamente 6;
  \task estén satisfechos al menos 6;
  \task estén satisfechos como mucho 5.
\end{tasks}

\ex{1} En una familia con 4 hijos, la probabilidad de que cada hijo sea chico es $0{,}5$. Calcula la probabilidad de que haya exactamente 2 chicos, y de que haya al menos una chica. ¿Cuántos chicos se esperan?

\ex{2} Una máquina fabrica un 3\,\% de piezas defectuosas. Las piezas se envasan en cajas de 50.\begin{tasks}(1)
  \task ¿Cuántas piezas defectuosas hay, de media, en una caja? ¿Con qué desviación típica?
  \task ¿Cuál es la probabilidad de que una caja no tenga ninguna pieza defectuosa? ¿Y de que tenga como mucho una?
\end{tasks}

\ex{3} ¿Cuántas veces, como mínimo, hay que lanzar una moneda para que la probabilidad de obtener al menos una cara sea mayor que $0{,}99$?


% =====================================================
\seccio{3}{Variables aleatorias continuas. Función de densidad}

\exf{1} Sea $f(x)=\frac16$ si $2\leq x\leq8$ (y $f(x)=0$ en otro caso).\begin{tasks}(1)
  \task Comprueba que es una función de densidad y dibújala.
  \task Calcula $P(X\lt4)$, $P(3\leq X\leq7)$ y $P(X=5)$.
  \task ¿Cuál es la media?
\end{tasks}\recorda{la probabilidad es el área bajo la función de densidad; el área total tiene que ser 1.}

\ex{2} Calcula $k$ para que $f(x)=k\,x$ si $0\leq x\leq4$ (y 0 en otro caso) sea una función de densidad. Después calcula $P(X\leq2)$ y $P(X\gt3)$.

\ex{1} El retraso de un tren (en minutos) se distribuye uniformemente entre 0 y 20 minutos. Calcula la probabilidad de que el retraso sea superior a 15 minutos, de que esté entre 5 y 12 minutos, y el retraso medio.

\ex{2} La función de densidad de una variable $X$ es un triángulo isósceles con base el intervalo $[0,4]$ y el vértice en $x=2$.\begin{tasks}(1)
  \task ¿Qué altura tiene que tener el triángulo?
  \task Calcula $P(X\leq1)$, $P(1\leq X\leq3)$ y $P(X\geq3{,}5)$.
\end{tasks}

\ex{2} Razona si las funciones siguientes pueden ser funciones de densidad (son 0 fuera del intervalo indicado):\begin{tasks}(1)
  \task $f(x)=0{,}5$ en $[0,3]$
  \task $f(x)=x-1$ en $[0,2]$
  \task $f(x)=0{,}25$ en $[1,5]$
\end{tasks}


% =====================================================
\seccio{4}{La distribución normal. La normal tipificada $N(0,1)$ y su tabla}

\exf{1} Si $Z\sim N(0,1)$, calcula con la tabla:\begin{tasks}(3)
  \task $P(Z\leq1{,}53)$
  \task $P(Z>2{,}05)$
  \task $P(Z<-0{,}82)$
  \task $P(Z\geq-1{,}2)$
  \task $P(0{,}75\leq Z\leq1{,}5)$
  \task $P(-2< Z<1)$
  \task $P(-2{,}58\leq Z\leq2{,}58)$
\end{tasks}\recorda{la tabla da $P(Z\leq z)$ para $z\geq0$. \ $P(Z\geq z)=1-P(Z\leq z)$, \ $P(Z\leq-z)=1-P(Z\leq z)$. Dibuja siempre la campana y el área que buscas.}

\ex{1} Comprueba con la tabla la regla 68--95--99,7: calcula $P(-1\leq Z\leq1)$, $P(-2\leq Z\leq2)$ y $P(-3\leq Z\leq3)$. ¿Qué significan estos resultados para una variable $X\sim N(\mu,\sigma)$?

\ex{2} Si $Z\sim N(0,1)$, halla el valor de $k$ en cada caso:\begin{tasks}(2)
  \task $P(Z\leq k)=0{,}8907$
  \task $P(Z\geq k)=0{,}0505$
  \task $P(Z\leq k)=0{,}25$
  \task $P(-k\leq Z\leq k)=0{,}9$
\end{tasks}\recorda{busca la probabilidad \emph{dentro} de la tabla; si es menor que 0,5, $k$ es negativo.}

\ex{2} Si $Z\sim N(0,1)$, calcula $P(Z\leq-1{,}5\ \text{o}\ Z\geq1{,}5)$, \ $P(-0{,}44\leq Z\leq0)$ \ y \ $P(Z\geq-2{,}35)$. Haz un dibujo de cada área.


% =====================================================
\seccio{5}{Tipificación. Problemas con la normal}

\exf{1} \pau{Illes Balears, PBAU extraordinaria 2025}Cuando un cliente entra en una tienda de electrodomésticos, el tiempo (en minutos) que pasa en ella sigue una distribución $N(15,10)$.\begin{tasks}(1)
  \task ¿Cuál es la probabilidad de que un cliente esté dentro de la tienda 25 minutos o menos?
  \task ¿Cuál es la probabilidad de que un cliente esté dentro de la tienda entre 5 y 25 minutos?
\end{tasks}\recorda{si $X\sim N(\mu,\sigma)$, entonces $Z=\dfrac{X-\mu}{\sigma}\sim N(0,1)$. Redondea $z$ a dos decimales y usa la tabla.}

\ex{1} \pau{Illes Balears, PBAU ordinaria 2021}El peso de las personas de un colegio mayor sigue una ley normal de media 70 kg y desviación típica 15 kg. Elegimos una persona al azar; calcula la probabilidad de que:\begin{tasks}(3)
  \task pese más de 80 kg;
  \task pese menos de 50 kg;
  \task pese entre 60 y 120 kg.
\end{tasks}

\ex{1} El coeficiente intelectual (CI) de los 2000 alumnos de una universidad sigue una distribución $N(100,15)$. ¿Cuántos alumnos se espera que tengan un CI superior a 130? ¿Y entre 85 y 115?

\ex{2} \pau{Illes Balears, PBAU extraordinaria 2023}En una población, la altura de los hombres sigue una normal de media 1,76 m y desviación típica 0,12 m, y la de las mujeres, una normal de media 1,62 m y desviación típica 0,11 m.\begin{tasks}(1)
  \task Elegimos un hombre al azar. ¿Cuál es la probabilidad de que su altura sea mayor o igual que 1,76 m?
  \task Elegimos una mujer al azar. ¿Cuál es la probabilidad de que su altura sea mayor o igual que 1,76 m?
  \task ¿Qué es más probable, que un hombre mida menos de 1,76 m o que lo haga una mujer?
\end{tasks}

\ex{2} \pau{Canarias, EBAU julio 2019}En una empresa hay 250 empleados. Su edad sigue una normal de media 44 años y desviación típica 18 años.\begin{tasks}(1)
  \task ¿Cuántos empleados se espera que tengan más de 62 años?
  \task ¿Cuántos empleados se espera que tengan menos de 40 años?
  \task ¿Cuántos empleados podrían conseguir el carnet joven de transporte si el requisito es ser mayor de edad y no haber cumplido los 30 años?
\end{tasks}

\ex{2} \pau{Illes Balears, PAU extraordinaria 2026}El tiempo que cada joven de una población dedica cada día a las redes sociales sigue una normal de media 116 minutos y desviación típica 42 minutos.\begin{tasks}(1)
  \task ¿Qué proporción de jóvenes dedica entre 100 y 200 minutos a las redes sociales?
  \task Juan es del 10\,\% de jóvenes que más tiempo les dedica. ¿Cuántos minutos les dedica, como mínimo?
\end{tasks}\recorda{en un problema inverso buscamos la probabilidad dentro de la tabla, leemos $z$ y deshacemos la tipificación: $x=\mu+z\,\sigma$.}

\ex{2} \pau{País Vasco, EAU ordinaria 2020}La altura de las mujeres de un país sigue una normal de media 163 cm y desviación típica 7 cm. Una empresa de disfraces quiere hacer cuatro tallas, cada una adecuada para el 25\,\% de las mujeres. ¿Qué alturas marcan el paso de una talla a otra?

\ex{2} \pau{País Vasco, EAU ordinaria 2020}El peso de las truchas de una piscifactoría sigue una normal de media 250 g y desviación típica 50 g. Solo son aptas para la venta las que superan un peso determinado.\begin{tasks}(1)
  \task ¿Cuál tiene que ser ese peso si se quiere que el 40\,\% de las truchas sean aptas para la venta?
  \task Si ese peso se establece en 280 g y en la piscifactoría hay 6000 truchas, ¿cuántas se podrán vender?
\end{tasks}

\ex{3} \pau{Castilla y León, EBAU julio 2018}En una empresa, el salario medio mensual es de 1400 euros y el 80,23\,\% de los trabajadores cobra menos de 1570 euros. Calcula la desviación típica del salario, suponiendo que sigue una distribución normal.

\ex{3} \pau{La Rioja, EBAU ordinaria 2023}La estatura de los policías de Francia sigue una normal de desviación típica 6,5 cm, de manera que el 11,507\,\% mide más de 183 cm.\begin{tasks}(1)
  \task Calcula la media de la estatura.
  \task Halla la estatura que es superada por el 88,493\,\% de los policías.
\end{tasks}


% =====================================================
\seccio{6}{Aproximación de la distribución binomial a la normal}

\exf{1} Sea $X\sim B(100;\ 0{,}2)$.\begin{tasks}(1)
  \task Comprueba que se puede aproximar por una normal y di cuál.
  \task Calcula aproximadamente $P(X\leq15)$, $P(X=20)$ y $P(X\geq25)$.
\end{tasks}\recorda{si $np\geq5$ y $nq\geq5$, $B(n,p)\approx N\left(np,\sqrt{npq}\right)$. No olvides la corrección de continuidad ($\pm0{,}5$).}

\ex{1} Lanzamos una moneda equilibrada 400 veces. Calcula la probabilidad de obtener menos de 190 caras.

\ex{2} Un examen tiene 100 preguntas de verdadero o falso. Para aprobar hay que acertar al menos 60. ¿Cuál es la probabilidad de aprobar contestando todas las preguntas al azar?

\ex{2} \pau{Andalucía, PAU ordinaria 2026 (adaptado)}Un experimento tiene éxito con probabilidad 0,4 y se repite 1350 veces de manera independiente. Sea $X$ el número de éxitos.\begin{tasks}(1)
  \task ¿Qué distribución sigue $X$?
  \task Calcula la probabilidad de que $X$ sea mayor que 499 y menor o igual que 580.
\end{tasks}

\ex{2} \pau{Canarias, EBAU ordinaria 2022 (adaptado)}Una compañía de seguros tiene asegurados 2500 coches. La probabilidad de que un coche tenga un accidente a lo largo de un año es 0,1. ¿Cuántos accidentes se esperan? ¿Cuál es la probabilidad de que el año que viene tengan un accidente al menos 270 coches?

\ex{3} Sea $X\sim B(50;\ 0{,}3)$. Calcula $P(10\leq X\leq20)$ con la aproximación normal. Con un ordenador se obtiene el valor exacto, $0{,}9120$. ¿Qué error se comete?

\vfill
\begin{center}\color{muted}\small\faCheckCircle[regular]\ \ ¡Ánimo! Comprueba tus resultados con la calculadora de probabilidades de los apuntes interactivos.\end{center}

\end{document}
